% Author: Akshat (condensed from report/section_data_tokenization_retrieval.md, parts T and C)
\section{Tokenization and candidate generation}
\label{sec:tok}

\subsection{How tokenizers split each script}

We measured the \emph{fertility}~\citep{rust2021good} of four multilingual tokenizers, i.e.\ subword tokens per
whitespace word, on val records grouped by country and script, and again on the Devanagari and
Malayalam records after romanisation (Table~\ref{tab:fertility}).

\begin{table}[h]
\centering\small
\resizebox{\linewidth}{!}{%
\begin{tabular}{lcccccc}
\toprule
 & US Latin & India Latin & Devanagari & Deva.\ romanised & Malayalam & Mal.\ romanised \\
\midrule
mBERT (WordPiece)           & 2.06 & 2.06 & \textbf{2.99} & 2.08 & \textbf{3.90} & 2.47 \\
XLM-R (SentencePiece)       & 2.05 & 1.94 & 2.17 & 2.06 & 2.82 & 2.36 \\
mDeBERTa-v3 (SentencePiece) & 1.95 & 1.87 & 2.45 & 2.15 & 2.81 & 2.41 \\
IndicBERTv2 (WordPiece)     & 2.04 & 1.77 & \textbf{1.54} & 1.92 & \textbf{1.71} & 2.32 \\
\bottomrule
\end{tabular}}
\caption{Tokens per word on val records.}
\label{tab:fertility}
\end{table}

mBERT fragments Indic scripts the most: 3.90 pieces per Malayalam word, falling back to single characters
and even splitting off the virama. IndicBERTv2 has the lowest fertility on both native scripts. For example,
it splits {\deva आदित्य प्रॉपर्टीज एलएलपी} (\emph{aditya properties llp}) into 5 pieces, against 10 for
XLM-R and 13 for mBERT. No tokenizer produces more than 0.2\% unknown tokens, so the cost of a foreign
script is fragmentation, not unknown words. Romanisation is a trade-off: it brings mBERT from 3.90 to 2.47
pieces per Malayalam word but makes IndicBERTv2 worse (1.71 to 2.32), because romanised Indian words are
in no tokenizer's vocabulary. This is why the cross-encoders get both the native name and its romanisation.

\subsection{Tokenization and normalisation for retrieval}

Table~\ref{tab:tfidf} compares TF-IDF analyzers and normalisation steps by the recall of the true S1 entity
among the top $k$ (same country, val).

\begin{table}[h]
\centering\small
\resizebox{\linewidth}{!}{%
\begin{tabular}{llcccr}
\toprule
analyzer & text & R@1 & R@5 & R@1 non-Latin & vocabulary \\
\midrule
words         & raw                           & 0.9729 & 0.9895 & 0.9536 & 59,722 \\
words         & romanised + abbrev.\ + lexicon & 0.9806 & 0.9930 & 0.9913 & 58,707 \\
char 3-grams  & romanised + abbrev.            & 0.9878 & 0.9960 & 0.9768 & 20,093 \\
\textbf{char 3-grams} & \textbf{romanised + abbrev.\ + lexicon} & \textbf{0.9910} & \textbf{0.9973} & \textbf{0.9952} & 20,086 \\
char 2--4-grams & romanised + abbrev.\ + lexicon & 0.9912 & 0.9972 & 0.9952 & 113,651 \\
XLM-R subwords & romanised + abbrev.\ + lexicon & 0.9894 & 0.9968 & 0.9931 & 15,403 \\
\bottomrule
\end{tabular}}
\caption{TF-IDF retrieval on val. ``Non-Latin'': S2/S3 records whose text is not in Latin script.}
\label{tab:tfidf}
\end{table}

The \textbf{transliteration lexicon} is learned from the training links. For each linked pair whose
romanised non-Latin name has as many words as its S1 name, the words are aligned by position. A romanised
word becomes an entry if it was aligned with the same English word at least 3 times and in at least 60\% of
its occurrences (\emph{praivtt}$\to$\emph{private}, \emph{knstrksns}$\to$\emph{constructions}). This is a
noisy-channel view~\citep{kernighan1990spelling}: a record is an S1 name transliterated to Latin letters,
then spelled phonetically. Romanisation inverts the first step, and the lexicon inverts the most common
errors of the second. On non-Latin records, the lexicon matters more than the tokenizer: it raises R@1 of
character 3-grams from 0.977 to 0.995, and with it all analyzers end within 0.004 of each other. Character
3-grams are the best analyzer overall. 2--4-grams tie at 5.7$\times$ the vocabulary, so we use 3-grams.

\subsection{Candidate generation}

Scoring all test S2/S3 records (about 100k) against all S1 entities of their country (about 11k) means
about $10^9$ pairs. Blocking keeps the 5 most likely S1 entities per record, and every matcher scores the
same 507,870 test pairs. We search only within the record's country, with three retrievers:
\begin{itemize}
  \item TF-IDF over character 3-grams of the normalised text;
  \item multilingual-e5-small~\citep{wang2024multilingual} (118M parameters), zero-shot;
  \item the same model fine-tuned for 3.5 minutes on 100k linked training pairs (part A), with a symmetric
        in-batch contrastive loss as in DPR~\citep{karpukhin2020dense}.
\end{itemize}
The TF-IDF and fine-tuned e5 lists are merged with reciprocal rank fusion~\citep{cormack2009reciprocal},
score $\sum 1/(60+\text{rank})$, and the top 5 are kept.

\begin{table}[h]
\centering\small
\begin{tabular}{lccc}
\toprule
retriever & R@1 & R@5 & R@10 \\
\midrule
TF-IDF char 3-gram   & 0.9912 & 0.9972 & 0.9983 \\
e5-small zero-shot   & 0.9883 & 0.9954 & 0.9968 \\
e5-small fine-tuned  & \textbf{0.9940} & 0.9988 & \textbf{0.9994} \\
RRF fusion, top 5    & -- & \textbf{0.9989} & -- \\
\bottomrule
\end{tabular}
\caption{Candidate retrieval on val.}
\label{tab:retrieval}
\end{table}

Zero-shot dense retrieval loses to character n-grams (R@1 0.9883 vs.\ 0.9912): a general embedding model has
not learned that two branches with similar names are different businesses. A short contrastive fine-tune
puts it ahead at every $k$ (Table~\ref{tab:retrieval}). \textbf{Recall is not the bottleneck}: a perfect
matcher on these candidates would reach macro \fhalf{} 0.99974 on val and \textbf{0.99957 on test}, so almost
all remaining error comes from scoring. About 14.7\% of candidate pairs are positive.
